\documentclass[11pt,a4paper]{article}
\usepackage[utf8]{inputenc}
\usepackage[T1]{fontenc}
\usepackage[english]{babel}
\usepackage{amsmath,amssymb,amsthm}
\usepackage{geometry}
\usepackage{titlesec}
\usepackage{fancyhdr}
\usepackage{microtype}
\usepackage{array}
\usepackage{booktabs}

\geometry{
  paperwidth=6.5in,
  paperheight=9.5in,
  textwidth=5in,
  textheight=7.8in,
  top=0.8in,
  bottom=0.8in,
  inner=0.75in,
  outer=0.75in,
}

\pagestyle{fancy}
\fancyhf{}
\fancyhead[LE,RO]{\thepage}
\fancyhead[CE]{\small A. CAYLEY---COLLECTED MATHEMATICAL PAPERS, VOL. III}
\fancyhead[CO]{\small\nouppercase{\leftmark}}
\renewcommand{\headrulewidth}{0pt}

\setcounter{page}{479}
\setcounter{section}{216}

\titleformat{\section}{\normalfont\large\bfseries\centering}{\thesection.}{0.5em}{}
\titleformat{\subsection}{\normalfont\normalsize\bfseries}{\thesubsection.}{0.5em}{}

\DeclareMathOperator{\cosec}{cosec}
\DeclareMathOperator{\VisViva}{Vis\,Viva}

\begin{document}
\thispagestyle{empty}

%%%%%%%%%%%%%%%%%%%%%%%%%%%%%%%%%%%%%%%%%%%%%%%%%%%%%%%%%%%%%%%%%%%%%%
% PAPER 217
%%%%%%%%%%%%%%%%%%%%%%%%%%%%%%%%%%%%%%%%%%%%%%%%%%%%%%%%%%%%%%%%%%%%%%
\section{A MEMOIR ON THE PROBLEM OF THE ROTATION OF A SOLID BODY}
\markboth{On the Rotation of a Solid Body.}{}

\begin{center}
[From the \textit{Philosophical Transactions}, vol.\ CLI.\ (1861), pp.~561--587.]
\end{center}

\subsection*{I.}

The theory of the rotation of a solid body about a fixed point may be considered as depending on the solution of the equations of motion in terms of certain elements; and for the complete development of the theory, these elements require to be connected with a set of angular coordinates adapted to geometrical and dynamical interpretation.

I commence by establishing the equations of motion under a form convenient for the subsequent treatment of the problem; and I then integrate these equations, first when the disturbing forces are neglected, and then when the disturbing forces are taken into account.

\subsection*{II.}

Denoting the coordinates, referred to axes fixed in space, by $x$, $y$, $z$, and the coordinates, referred to axes fixed in the body, by $x_1$, $y_1$, $z_1$, the relation between the two sets of coordinates is expressed by means of a table of coefficients:
\[
\begin{array}{c|ccc}
 & x & y & z \\
\hline
x_1 & \alpha & \beta & \gamma \\
y_1 & \alpha' & \beta' & \gamma' \\
z_1 & \alpha'' & \beta'' & \gamma''
\end{array}
\]
so that
\[
x_1 = \alpha x + \beta y + \gamma z, \quad
y_1 = \alpha' x + \beta' y + \gamma' z, \quad
z_1 = \alpha'' x + \beta'' y + \gamma'' z,
\]
and in like manner
\[
x = \alpha x_1 + \alpha' y_1 + \alpha'' z_1, \quad
y = \beta x_1 + \beta' y_1 + \beta'' z_1, \quad
z = \gamma x_1 + \gamma' y_1 + \gamma'' z_1.
\]

If $p$, $q$, $r$ denote the components of angular velocity about the axes fixed in the body, and if $X$, $Y$, $Z$ denote the moments of the impressed forces about these axes respectively, then the equations of motion are
\begin{align*}
A\frac{dp}{dt} + (C-B)qr &= X, \\
B\frac{dq}{dt} + (A-C)rp &= Y, \\
C\frac{dr}{dt} + (B-A)pq &= Z,
\end{align*}
where $A$, $B$, $C$ are the principal moments of inertia at the fixed point.

The connection between the differential coefficients of the coefficients $\alpha$, $\beta$, etc., and the quantities $p$, $q$, $r$ is given by the following formulae; writing for shortness $(1)$, $(2)$, $(3)$ to denote $x_1$, $y_1$, $z_1$ and $(1')$, $(2')$, $(3')$ to denote the differential coefficients of these quantities in regard to $t$:
\begin{align*}
(1') &= q(3) - r(2), \\
(2') &= r(1) - p(3), \\
(3') &= p(2) - q(1),
\end{align*}
from which, attending to the values of $(1)$, $(2)$, $(3)$ in terms of $x$, $y$, $z$, we may obtain the differential coefficients of the coefficients $\alpha$, $\beta$, etc.

The position of the body may be determined by means of the position of the principal axes, or say by means of the position of the axes $(x_1, y_1, z_1)$; and this is most conveniently determined by means of the spherical triangle formed by these axes with a fixed plane and a fixed line in that plane, or (as in the theory of the top) by means of the angles $\theta$, $\phi$, $\psi$, where $\theta$ is the inclination of the plane of $(x_1, y_1)$ to the fixed plane, $\phi$ is the longitude of the node of the plane of $(x_1, y_1)$ upon the fixed plane, and $\psi$ is the angle which the axis of $x_1$ makes with the node.

But it is more convenient to introduce other angular coordinates, and in fact to consider the problem as analogous to the problem of elliptic motion.

Suppose that the body is rotating about the axis of $z_1$ in the positive direction; and let $h$, $k$ denote respectively the constants of Vis Viva and of areas; so that
\[
Ap^2 + Bq^2 + Cr^2 = h, \qquad A^2p^2 + B^2q^2 + C^2r^2 = k^2.
\]

The plane perpendicular to the axis of resultant angular momentum is called the \textit{principal plane}; and the intersection of this plane with the plane of $(x_1, y_1)$ is a line which may be called the \textit{line of nodes}.

I introduce the following angular coordinates:
\begin{itemize}
\item[$T$,] the departure of the line of nodes in the plane of $(x_1, y_1)$;
\item[$S$,] the inclination of the plane of $(x_1, y_1)$ to the principal plane;
\item[$F$,] the departure of the axis of $x_1$ from the line of nodes;
\item[$\theta$,] the longitude of the node of the principal plane on the fixed plane;
\item[$\phi$,] the inclination of the principal plane to the fixed plane;
\item[$\sigma$,] the departure of the node of the principal plane.
\end{itemize}

The position of the body is determined by $T$, $S$, $F$, together with the position of the principal plane determined by $\theta$, $\phi$, $\sigma$.

And if $A$, $B$, $C$ are the moments of inertia, then the Vis Viva function or sum of the elements of mass, each into the half square of its velocity, is
\[
T = \tfrac{1}{2}(Ap^2 + Bq^2 + Cr^2);
\]
and if $\Omega$ be the disturbing function (taken with Lagrange's sign), then the equations of motion
\[
\frac{d}{dt}\frac{dT}{d\dot{q}} - \frac{dT}{dq} = \frac{d\Omega}{dq}
\]
give as usual
\begin{multline}
A\frac{dp}{dt} + (C-B)qr = -\sin S\left(\cosec F\,\frac{d\Omega}{dT} + \cot F\,\frac{d\Omega}{dS}\right) \notag\\
+ \cos S\,\frac{d\Omega}{dF},
\end{multline}
which, with the equations for $p$, $q$, $r$, determine the motion of the body.

\subsection*{III.}

First, to integrate the equations of motion when the disturbing forces are neglected; we have, as usual, the integral equations
\[
Ap^2 + Bq^2 + Cr^2 = h, \qquad A^2p^2 + B^2q^2 + C^2r^2 = k^2;
\]
where $h$, $k$ are constants of integration, viz., $h$ is the constant of Vis Viva, and $k$ is the constant of the principal area.

Moreover, if the coefficients $\alpha$, $\beta$, etc.\ are those which belong to the transformation from $xyz$ to $x_1 y_1 z_1$, viz., if in the table
\[
\begin{array}{c|ccc}
 & x & y & z \\
\hline
x_1 & \alpha & \beta & \gamma \\
y_1 & \alpha' & \beta' & \gamma' \\
z_1 & \alpha'' & \beta'' & \gamma''
\end{array}
\]
the values of the coefficients are
\begin{align*}
\alpha &= \cos T, & \beta &= \sin T, & \gamma &= 0, \\
\alpha' &= -\sin T\cos S, & \beta' &= \cos T\cos S, & \gamma' &= \sin S, \\
\alpha'' &= \sin T\sin S, & \beta'' &= -\cos T\sin S, & \gamma'' &= \cos S,
\end{align*}
so that the relation between the coordinates $xyz$ and $x_1 y_1 z_1$ is given by the table
\[
\begin{array}{c|ccc}
 & x_1 & y_1 & z_1 \\
\hline
x & \cos T & -\sin T\cos S & \sin T\sin S \\
y & \sin T & \cos T\cos S & -\cos T\sin S \\
z & 0 & \sin S & \cos S
\end{array}
\]
where the values of the coefficients are given by
\begin{align*}
\alpha &= \cos T, & \alpha' &= -\sin T\cos S, & \alpha'' &= \sin T\sin S, \\
\beta &= \sin T, & \beta' &= \cos T\cos S, & \beta'' &= -\cos T\sin S, \\
\gamma &= 0, & \gamma' &= \sin S, & \gamma'' &= \cos S.
\end{align*}

The two equations
\[
Ap^2 + Bq^2 + Cr^2 = h, \qquad A^2p^2 + B^2q^2 + C^2r^2 = k^2,
\]
may be considered as determining $p$, $q$, in terms of $r$, and the equation
\[
C\frac{dr}{dt} + (B-A)pq = 0,
\]
then gives
\[
\frac{dr}{dt} = -\frac{(B-A)pq}{C};
\]
whence also the equation for $dT_2$ becomes
\[
k\,dT_2 = \frac{k(h-Cr^2) - C^2r^2}{(B-A)pq}\,dr.
\]

Instead of the time $t$, I consider a function $g = nt + \text{const.}$, $n$ being for the present an arbitrary constant quantity which may be a function of the constants $h$ and $k$; we have thus
\[
dg = n\,dt.
\]

And the integral equation is
\[
g = n\int\frac{-C\,dr}{(B-A)pq},
\]
where it is assumed that the integrals are each of them taken from $r = r_0$, $r_0$ being an arbitrary constant value, say a function of $h$ and $k$.

The quantity $g$ is analogous to the mean anomaly in elliptic motion, or rather it will become so when the significations of $n$ and $r_0$ are fixed; it is considered as implicitly involving a constant of integration, and, consequently, no constant of integration is added to the integral: as regards $T_2$, the constant of integration is $\sigma$, which denotes the initial value (corresponding to $r = r_0$) of the angle $T_2$.

\subsection*{IV.}

Recapitulating the integral equations, we have
\begin{align*}
Ap^2 + Bq^2 + Cr^2 &= h, \\
A^2p^2 + B^2q^2 + C^2r^2 &= k^2, \\
Ap &= k\sin S_2\sin F_2, \\
Bq &= k\cos S_2\sin F_2, \\
Cr &= k\cos F_2, \\
g &= n\int\frac{-C\,dr}{(B-A)pq},
\end{align*}
which equations give $r$, $p$, $q$, and thence $S_2$ and $F_2$, in terms of $g$ and the constants $h$ and $k$.

Moreover,
\[
T = \sigma + \int\frac{k(h-Cr^2)-C^2r^2}{k(B-A)pq}\,dr,
\]
which gives $T$ in terms of $g$ and the constants $h$, $k$, $\sigma$; and then $T$, $S$, $F$ are given in terms of $T_2$, $S_2$, $F_2$, and the constants $\theta$, $\sigma$, $\phi$, as follows, viz., we have a spherical triangle $ABC$, the sides whereof $a$, $b$, $c$, and the opposite angles $A$, $B$, $C$, are respectively,
\begin{center}
\begin{tabular}{ll}
\textit{Sides} & \textit{Opposite angles} \\
$\sigma - S$, & $\phi$, \\
$T_2 - \sigma$, & $180^\circ - \psi$, \\
$T - \theta$, & $F$,
\end{tabular}
\end{center}
as appears by the figure.

The above values of $p$, $q$, $r$, give
\[
k^2 = A^2p^2 + B^2q^2 + C^2r^2 = k^2\sin^2 F_2 + C^2r^2,
\]
an equation which will be useful in the sequel.

The coefficient $n$, and the value $r_0$, which is the inferior limit of the integrals, are thus far considered as arbitrary functions of $h$ and $k$. As already remarked, $g$ (which is a variable quantity, $= nt + c$) is used as an element, and the elements are $h$, $k$, $g$, $\sigma$, $\theta$, $\phi$. As to the signification of the different elements, it is proper to remark that it is not for the purposes of the present memoir necessary to enter into the complete theory of the rotation of a solid body; the problem here considered is that of the variation of the elements for a solid body rotating about a fixed point, and acted on by a disturbing function.

\subsection*{V.}

This being premised, we have (corresponding to the orbit in the theory of elliptic motion) the principal plane, with a departure point therein, the positions whereof, in respect to the fixed plane of reference, are determined by $\theta$, the longitude of node; $\phi$, the inclination; and $\sigma$, the departure of node.

The precise signification of $\sigma$ depends upon that of $r_0$, and the signification of $r_0$ being assumed to be completely determined, that of $\sigma$ will be so likewise; $\sigma$ is to be considered as an angle measured in the principal plane from the departure point, and determining in that plane a line (or, treating the plane as an orbit, a point) which I call the \textit{rotation pericentre}, or simply the \textit{pericentre}; say $\sigma$ is the departure of the pericentre; and then $g$ is an angle varying uniformly with the time, used for expressing in terms of the time the angle $T$, which determines the position, in regard to the principal plane and departure point therein, of the node of the plane of $x_1 y_1$, or equator, upon the principal plane, such node corresponding with the moving body in the theory of elliptic motion.

We may in fact say: $T$, the departure of the last-mentioned node, $= \sigma$, the departure of pericentre, $+ (T-\sigma)$, the rotation true anomaly of such node; $T_2$ being a function of $g$, the rotation mean anomaly of such node.

The elements are then as follows, viz.:
\begin{itemize}
\item[$h$,] the constant of Vis Viva,
\item[$k$,] the constant of principal area,
\item[$\theta$,] the longitude of node of the principal plane,
\item[$\phi$,] the inclination of the principal plane,
\item[$\sigma$,] the departure of the pericentre,
\item[$g$,] the mean anomaly.
\end{itemize}

Moreover, corresponding to the theory of elliptic motion, we have also the anomalies $T_2$, $S_2$, $F_2$, and the pericentre $S$, $F$, viz.:
\begin{itemize}
\item[$T_2$,] the true anomaly of the node on the equator,
\item[$S_2$,] the inclination of the equator to the principal plane,
\item[$F_2$,] the departure of the axis $x_1$ from the node,
\item[$S$,] the inclination of the equator (constant),
\item[$F$,] the departure of the axis $x_1$ from the node (constant).
\end{itemize}

We have $S$, $F$ as constants determined by the initial values; and $T_2$, $S_2$, $F_2$ are functions of $g$ (and therefore of $t$) determined by the equations of motion.

\subsection*{VI.}

If, now, the Disturbing Function is taken into account, the equations are to be integrated by the method of the variation of the elements.

I use $\delta g$ to denote that part of the variation of $g$ ($= t\,\delta n + \delta c$, if $g = nt + c$), which depends on the variation of the constants, and in like manner for $\delta p$, $\delta q$, $\delta r$, etc.

I assume, moreover, that $\sigma$ is a departure point in the principal plane; this gives $d\sigma - \cos\phi\,d\theta = 0$; and we see that the equations which lead to the expressions for the variations of the elements are
\begin{align*}
A\,dp &= -\sin S\left(\cosec F\,\frac{d\Omega}{dT} + \cot F\,\frac{d\Omega}{dS}\right)dt + \cos S\,\frac{d\Omega}{dF}\,dt, \\
B\,dq &= \cos S\left(\cosec F\,\frac{d\Omega}{dT} + \cot F\,\frac{d\Omega}{dS}\right)dt + \sin S\,\frac{d\Omega}{dF}\,dt, \\
C\,dr &= -\frac{d\Omega}{dF}\,dt, \\
dT &= 0, \quad dS = 0, \quad dF = 0, \quad d\sigma - \cos\phi\,d\theta = 0,
\end{align*}
where, of course, $\Omega = \Omega(T, S, F)$.

The variations $dh$ and $dk$ are obtained from the equations
\[
dh = Ap\,dp + Bq\,dq + Cr\,dr, \qquad k\,dk = A^2p\,dp + B^2q\,dq + C^2r\,dr,
\]
expressions which will presently be transformed.

\subsection*{VII.}

Proceeding to carry out the foregoing plan, the equation
\[
dh = Ap\,dp + Bq\,dq + Cr\,dr,
\]
putting for $A\,dp$, $B\,dq$, $C\,dr$, their values, gives
\begin{multline}
dh = (p\sin S + q\cos S)\left(\cosec F\,\frac{d\Omega}{dT} + \cot F\,\frac{d\Omega}{dS}\right)dt \\
+ (p\cos S + q\sin S)\,\frac{d\Omega}{dF}\,dt - r\,\frac{d\Omega}{dF}\,dt,
\end{multline}
which may also be written
\begin{multline}
dh = \cos(S_2-S)\sin F_2\left(\cosec F\,\frac{d\Omega}{dT} + \cot F\,\frac{d\Omega}{dS}\right)dt \\
+ \sin(S_2-S)\sin F_2\,\frac{d\Omega}{dF}\,dt - \cos F_2\,\frac{d\Omega}{dF}\,dt.
\end{multline}

In fact, if, in this expression, we consider first the terms involving $\dfrac{d\Omega}{dF}\,dt$, the coefficient is
\[
-Cr\frac{(h-Cr^2)}{k(B-A)pq} + \cdots
\]
and the quantities in $\{\}$ being respectively $p(k^2-C^2r^2)$ and $q(k^2-C^2r^2)$, and since $k^2 - C^2r^2 = k^2\sin^2 F_2$, the coefficient is $= p\cos S + q\sin S$.

In like manner, for the terms involving $\left(\cosec F\,\dfrac{d\Omega}{dT} + \cot F\,\dfrac{d\Omega}{dS}\right)dt$, the coefficient is
\begin{multline*}
\frac{1}{k\sin F_2}\bigl[(h-Cr^2)\cos(S_2-S) + (B-A)pq\sin(S_2-S)\bigr] \\
= \frac{1}{k\sin F_2}\bigl[(h-Cr^2)\cos S_2\cos S + (h-Cr^2)\sin S_2\sin S \\
+ (B-A)pq(\sin S_2\cos S - \cos S_2\sin S)\bigr];
\end{multline*}
which is $= p\sin S + q\cos S$; and the foregoing transformed expression for $dh$ is thus seen to be correct.

The equation
\[
k\,dk = A^2p\,dp + B^2q\,dq + C^2r\,dr,
\]
substituting for $A\,dp$, $B\,dq$, $C\,dr$, their values, becomes
\begin{multline*}
k\,dk = Ap\bigl\{\sin S(\cosec F\,dT + \cot F\,dS) + \cos S\,dF\bigr\} \\
+ Bq\bigl\{\cos S(\cosec F\,dT + \cot F\,dS) + \sin S\,dF\bigr\} - Cr\,dF.
\end{multline*}

But, from the equations
\[
Ap = k\sin S_2\sin F_2, \quad Bq = k\cos S_2\sin F_2, \quad Cr = k\cos F_2,
\]
we deduce
\begin{align*}
Ap\sin S + Bq\cos S &= k\cos(S_2-S)\sin F_2, \\
Ap\cos S + Bq\sin S &= k\sin(S_2-S)\sin F_2, \\
Cr &= k\cos F_2,
\end{align*}
and we thence find
\begin{multline}
dk = \cos(S_2-S)\sin F_2\left(\cosec F\,\frac{d\Omega}{dT} + \cot F\,\frac{d\Omega}{dS}\right)dt \notag\\
+ \sin(S_2-S)\sin F_2\,\frac{d\Omega}{dF}\,dt - \cos F_2\,\frac{d\Omega}{dF}\,dt,
\end{multline}
so that $dh$ and $dk$ are now determined.

\subsection*{VIII.}

The expressions for $dh$ and $dk$ being thus obtained, and $dr$ being also given by the equation $C\,dr = -\dfrac{d\Omega}{dF}\,dt$, we may now express $dF_2$ and $dS_2$ in terms of $dh$, $dk$, $dr$.

In fact, the equations
\[
Cr = k\cos F_2, \qquad \frac{Cr}{k} = \cos F_2,
\]
give immediately
\[
k\sin F_2\,dF_2 = -C\,dr + \cos F_2\,dk,
\]
or, multiplying by $k^2 - C^2r^2$ ($= k^2\sin^2 F_2$), and replacing $\cos S_2\sin F_2$, $\sin S_2\sin F_2$ by their values $-ABpq$, and dividing, the last equation becomes
\[
dF_2 = \frac{-C(h-Cr^2)r\,dr + C^2r\,dr - C(h-Cr^2)\,dk + \cdots}{k(k^2-C^2r^2)(B-A)pq},
\]
and thence also
\[
\cos F_2\,dS_2 = \frac{-C^2(h-Cr^2)r^2\,dr + \cdots}{k(k^2-C^2r^2)(B-A)pq},
\]
which is the value of $dT_2$, as given by the equation (not yet demonstrated) $dT_2 = \cos F_2\,dF_2$.

The expression for $dF_2$, substituting for $dr$ and $dk$, their values, gives
\begin{multline*}
k\sin F_2\,dF_2 = \cos F_2\sin F_2\cos(S_2-S)\left(\cosec F\,\frac{d\Omega}{dT} + \cot F\,\frac{d\Omega}{dS}\right)dt \\
- \cos F_2\sin F_2\sin(S_2-S)\,\frac{d\Omega}{dF}\,dt + \cdots
\end{multline*}
or, combining the last two terms, and reducing,
\begin{multline*}
dF_2 = \frac{1}{k}\bigl\{\cos F_2\cos(S_2-S)(\cosec F\,d\Omega/dT + \cot F\,d\Omega/dS)\,dt \\
- \cos F_2\sin(S_2-S)\,d\Omega/dF\,dt\bigr\}.
\end{multline*}

And the value of the numerator being
\begin{multline*}
k\sin S_2\sin F_2\bigl\{-\cos S(\cosec F\,dT + \cot F\,dS) + \sin S\,dF\bigr\} \\
+ k\cos S_2\sin F_2\bigl\{-\sin S(\cosec F\,dT + \cot F\,dS) - \cos S\,dF\bigr\},
\end{multline*}
we find
\begin{multline*}
dS_2 = \frac{1}{k}\bigl\{-\sin(S_2-S)(\cosec F\,d\Omega/dT + \cot F\,d\Omega/dS)\,dt \\
+ \cos(S_2-S)\,d\Omega/dF\,dt\bigr\};
\end{multline*}
the last-mentioned expressions for $dT_2$ and $dS_2$ will be used for obtaining $d\theta$, $d\sigma$, $d\phi$.

\subsection*{IX.}

I form now the equations
\begin{align*}
-\sin S\sin F\,dT + \cos S\,dF &= p\,dt - \sin S_2\sin F_2\,dT_2 + \cos S_2\,dF_2 \\
&\quad + \alpha_2(\sin\sigma\sin\phi\,d\theta + \cos\sigma\,d\phi) \\
&\quad + \alpha'_2(\cos\sigma\sin\phi\,d\theta + \sin\sigma\,d\phi) \\
&\quad + \alpha''_2(-d\sigma + \cos\phi\,d\theta), \\
\cos S\sin F\,dT + \sin S\,dF &= q\,dt + \cos S_2\sin F_2\,dT_2 + \sin S_2\,dF_2 \\
&\quad + \beta_2(\sin\sigma\sin\phi\,d\theta + \cos\sigma\,d\phi) \\
&\quad + \beta'_2(\cos\sigma\sin\phi\,d\theta + \sin\sigma\,d\phi) \\
&\quad + \beta''_2(-d\sigma + \cos\phi\,d\theta), \\
-dS + \cos F\,dT &= r\,dt - dS_2 + \cos F_2\,dT_2 \\
&\quad + \gamma_2(\sin\sigma\sin\phi\,d\theta + \cos\sigma\,d\phi) \\
&\quad + \gamma'_2(\cos\sigma\sin\phi\,d\theta + \sin\sigma\,d\phi) \\
&\quad + \gamma''_2(-d\sigma + \cos\phi\,d\theta).
\end{align*}

For the present purpose we are to omit the terms $p\,dt$, $q\,dt$, $r\,dt$, and to write $dT = 0$, $dS = 0$, $dF = 0$.

We have thus
\begin{align*}
- (\sin S_2\sin F_2\,dT_2 + \cos S_2\,dF_2) &= \alpha_2(-\sin\sigma\sin\phi\,d\theta + \cos\sigma\,d\phi) \\
&\quad + \alpha'_2(\cos\sigma\sin\phi\,d\theta + \sin\sigma\,d\phi) \\
&\quad + \alpha''_2(-d\sigma + \cos\phi\,d\theta), \\
- (\cos S_2\sin F_2\,dT_2 + \sin S_2\,dF_2) &= \beta_2(-\sin\sigma\sin\phi\,d\theta + \cos\sigma\,d\phi) \\
&\quad + \beta'_2(\cos\sigma\sin\phi\,d\theta + \sin\sigma\,d\phi) \\
&\quad + \beta''_2(-d\sigma + \cos\phi\,d\theta), \\
-(-dS_2 + \cos F_2\,dT_2) &= \gamma_2(-\sin\sigma\sin\phi\,d\theta + \cos\sigma\,d\phi) \\
&\quad + \gamma'_2(\cos\sigma\sin\phi\,d\theta + \sin\sigma\,d\phi) \\
&\quad + \gamma''_2(-d\sigma + \cos\phi\,d\theta).
\end{align*}

Hence we have
\begin{multline*}
\sin\sigma\sin\phi\,d\theta + \cos\sigma\,d\phi = \alpha_2(\sin S_2\sin F_2\,dT_2 + \cos S_2\,dF_2) \\
- \beta_2(\cos S_2\sin F_2\,dT_2 + \sin S_2\,dF_2) - \gamma_2(\cos F_2\,dT_2 - dS_2),
\end{multline*}
which, attending to the values of $\alpha''_2$, $\beta''_2$, $\gamma''_2$, may be written,
\begin{multline*}
\sin\sigma\sin\phi\,d\theta + \cos\sigma\,d\phi = \alpha_2(\alpha''_2\,dT_2 + \cos S_2\,dF_2) \\
- \beta_2(\beta''_2\,dT_2 + \sin S_2\,dF_2) - \gamma_2(\gamma''_2\,dT_2 - dS_2) \\
= -(\alpha_2\cos S_2 + \beta_2\sin S_2)\,dF_2 + \gamma_2\,dS_2,
\end{multline*}
since the coefficient of $dT_2$ vanishes; and, in like manner,
\begin{align*}
\cos\sigma\sin\phi\,d\theta + \sin\sigma\,d\phi &= -(\alpha'_2\cos S_2 + \beta'_2\sin S_2)\,dF_2 + \gamma'_2\,dS_2, \\
-d\sigma + \cos\phi\,d\theta &= -dT_2 - (\alpha''_2\cos S_2 + \beta''_2\sin S_2)\,dF_2 + \gamma''_2\,dS_2.
\end{align*}

But we have
\begin{align*}
\alpha_2\cos S_2 + \beta_2\sin S_2 &= \cos T_2\cos S_2 + \sin T_2\cos S_2 = \cdots \\
\alpha'_2\cos S_2 + \beta'_2\sin S_2 &= -\sin T_2\cos S_2 + \cos T_2\cos S_2 = \cdots
\end{align*}

\subsection*{X.}

The foregoing formulae become more simple when $n$ is a function of $h$ only, and when the relation between $r$ and $g$ is that which belongs to elliptic motion; that is, when
\[
\frac{dg}{dt} = n = \sqrt{\frac{2h-2\Omega}{k^2}},
\]
where $\Omega$ is the disturbing function. In this case we have
\[
\frac{d^2g}{dt^2} = 2n\frac{d\Omega}{dh},
\]
and the formula for $dg$ reduces to
\[
dg = 2n\frac{d\Omega}{dh}\,dt.
\]

\subsection*{XI.}

In the original expression for $k\sin\phi\,d\theta$, the three coefficients in $[\,]$ are
\begin{align*}
\cos b\sin a - \sin b\cos a\cos C &= \cos B\sin c = -\cos F\sin(T-\theta), \\
\cos b\cos a + \sin b\sin a\cos C &= \cos c = \cos(T-\theta), \\
\sin b\sin C &= \sin B\sin c = -\sin T\sin(T-\theta),
\end{align*}
and thence
\begin{align*}
k\sin\phi\,d\theta &= \cos F\sin(T-\theta)\left(\cosec F\,\frac{d\Omega}{dT} + \cot F\,\frac{d\Omega}{dS}\right)dt - \cos(T-\theta)\,\frac{d\Omega}{dF}\,dt \\
&= \sin(T-\theta)\,dT - \cosec F\sin(T-\theta)\,\frac{d\Omega}{dt}\,dt,
\end{align*}
or we have finally
\[
k\sin\phi\,d\theta = -\sin(T-\theta)\left(\cot F\,\frac{d\Omega}{dS} + \cosec F\,\frac{d\Omega}{dT}\right)dt.
\]

The expression for $d\sigma$ can be obtained from either of those for $d\theta$, by the equation $d\sigma = \cos\phi\,d\theta$.

\subsection*{XII.}

The comparison of the terms involving $dm$ gives at once
\[
\frac{d\Omega}{dm} = X,
\]
or, substituting for $C$ its value,
\begin{multline*}
\frac{d\Omega}{dm} = \cos(S_2-S)\sin F_2\left(\cosec F\,\frac{d\Omega}{dT} + \cot F\,\frac{d\Omega}{dS}\right) \\
- \sin(S_2-S)\sin F_2\,\frac{d\Omega}{dF} - \cos F_2\,\frac{d\Omega}{dS};
\end{multline*}
and comparing with the expression for $d\sigma$, we find
\[
\frac{d\Omega}{dk} = -\frac{1}{k}\frac{d\Omega}{dt}.
\]

The terms involving $dr$ give
\[
\frac{d\Omega}{dr} = \frac{k(B-A)pq}{k^2-C^2r^2}\left(\frac{d\Omega}{dh} + \frac{n}{k}\frac{d\Omega}{dk}\right) - \frac{2n}{k(B-A)pq}\frac{d\Omega}{dS},
\]
or, substituting for $X$, $Y$, $Z$, their values,
\begin{multline*}
\frac{d\Omega}{dr} = \frac{k(B-A)pq}{k^2-C^2r^2}\left(\frac{d\Omega}{dh} + \frac{n}{k}\frac{d\Omega}{dk}\right)\sin(S_2-S)\sin F_2\left(\cosec F\,\frac{d\Omega}{dT} + \cot F\,\frac{d\Omega}{dS}\right) \\
+ \cos(S_2-S)\sin F_2\,\frac{d\Omega}{dF} - \frac{2n}{k(B-A)pq}\frac{d\Omega}{dS},
\end{multline*}
which agrees with the second form for $\dfrac{d\Omega}{dk}$, and gives as before
\[
\frac{dg}{dt} = n.
\]

The terms involving $dh$ give
\[
dh = 2n\,\frac{d\Omega}{dg}\,dt;
\]
and thence
\[
\frac{dg}{dt} = n - \frac{Z}{k(B-A)pq},
\]
or, substituting for $2n$ its value $\dfrac{dg}{dt}$, for $X$ the value $\dfrac{d\Omega}{dm}$, and for $Z$ its value $\dfrac{d\Omega}{dr}$, we have
\[
\frac{dg}{dt} = n - \frac{1}{k(B-A)pq}\frac{d\Omega}{dr}.
\]

\subsection*{XIII.}

We have
\[
\frac{d^2r}{dt^2} = \frac{k^2(h-Cr^2)-C^2r^2}{(B-A)^2p^2q^2}\cdot\frac{dr}{dt};
\]
whence, observing that
\[
\frac{dr}{dt} = -\frac{(B-A)pq}{C},
\]
we obtain
\[
\frac{d^2r}{dt^2} = \frac{k^2(h-Cr^2)-C^2r^2}{C(B-A)pq}.
\]

\subsection*{XIV.}

To verify the formula for $dg$, we have
\[
dg = \frac{k(B-A)pq\left(\dfrac{d\Omega}{dh} + \dfrac{n}{k}\dfrac{d\Omega}{dk}\right) - 2n\dfrac{d\Omega}{dS}}{k(B-A)pq}\,dt.
\]

Now
\[
\frac{d\Omega}{dh} = \frac{1}{2n}\frac{dg}{dt}, \qquad \frac{d\Omega}{dk} = -\frac{1}{k}\frac{d\Omega}{dt},
\]
and substituting these values, the formula becomes
\[
dg = \left(\frac{dg}{dt} - \frac{2n}{k(B-A)pq}\frac{d\Omega}{dS}\right)dt,
\]
which is the required result.

\subsection*{XV.}

The final results for the variations of the elements may be collected as follows:
\begin{align}
dh &= 2n\,\frac{d\Omega}{dg}\,dt, \\
dk &= \cos(S_2-S)\sin F_2\left(\cosec F\,\frac{d\Omega}{dT} + \cot F\,\frac{d\Omega}{dS}\right)dt \notag\\
&\quad - \sin(S_2-S)\sin F_2\,\frac{d\Omega}{dF}\,dt - \cos F_2\,\frac{d\Omega}{dS}\,dt, \\
dg &= \left(\frac{d\Omega}{dh} + \frac{n}{k}\frac{d\Omega}{dk}\right)dt - \frac{2n}{k(B-A)pq}\frac{d\Omega}{dS}\,dt \notag\\
&\quad + \frac{2n}{k(B-A)pq}\bigl\{\sin(S_2-S)\sin F_2(\cosec F\,d\Omega/dT \notag\\
&\quad\qquad + \cot F\,d\Omega/dS) + \cos(S_2-S)\sin F_2\,d\Omega/dF\bigr\}\,dt, \\
d\psi &= \frac{1}{k\sin\phi}\frac{d\Omega}{d\theta}\,dt, \\
d\theta &= -\frac{1}{k\sin\phi}\left(\cot F\,\frac{d\Omega}{dS} + \cosec F\,\frac{d\Omega}{dT}\right)dt, \\
d\sigma &= \frac{\cos\phi}{k\sin\phi}\left(\cot F\,\frac{d\Omega}{dS} + \cosec F\,\frac{d\Omega}{dT}\right)dt, \\
d\phi &= \frac{1}{k}\frac{d\Omega}{d\psi}\,dt.
\end{align}

These formulae express the variations of the elements in terms of the differential coefficients of the disturbing function with respect to the elements themselves; and they constitute the complete solution of the problem of the rotation of a solid body, when the disturbing function is given.

\vfill
\pagebreak

%%%%%%%%%%%%%%%%%%%%%%%%%%%%%%%%%%%%%%%%%%%%%%%%%%%%%%%%%%%%%%%%%%%%%%
% PAPER 218
%%%%%%%%%%%%%%%%%%%%%%%%%%%%%%%%%%%%%%%%%%%%%%%%%%%%%%%%%%%%%%%%%%%%%%
\setcounter{section}{217}
\section{A THIRD MEMOIR ON THE PROBLEM OF DISTURBED ELLIPTIC MOTION}
\markboth{On Disturbed Elliptic Motion.}{}

\begin{center}
[From the \textit{Philosophical Transactions}, vol.\ CLI.\ (1861), pp.~589--610.]
\end{center}

\subsection*{I.}

The position in respect to a fixed plane of reference and origin of longitudes therein, of the variable ecliptic and of the departure point or origin of longitudes therein, are determined by $\theta$, the longitude of node, $\sigma$, the departure of node, $\phi$, the inclination, where, by the definition of a departure point,
\[
d\sigma - \cos\phi\,d\theta = 0,
\]
and then, in respect to the variable ecliptic and departure point or origin of longitudes therein, the position of the disturbed body is determined by $r$, the radius vector, $v$, the longitude, $\gamma$, the latitude; and this being so, then (Supplementary Memoir, pp.~220, 227) the expression for the Vis Viva function is,
\[
2T = \dot{r}^2 + r^2\dot{\gamma}^2 + r^2\cos^2\gamma\,\dot{v}^2 + 2Q\dot{r} + 2Rr\dot{v},
\]
where
\begin{align*}
Q &= -\dot{\gamma} + [\cos(v-\sigma)\sin\phi\cdot\dot{\theta} - \sin(v-\sigma)\dot{\phi}], \\
R &= \cos\gamma\cdot\dot{v} - \sin\gamma[\sin(v-\sigma)\sin\phi\cdot\dot{\theta} + \cos(v-\sigma)\dot{\phi}],
\end{align*}
the superscript dots being used to denote differentiation with respect to the time.

The last-mentioned expressions may for shortness be denoted by
\[
Q = -\dot{\gamma} + A\sin\gamma, \qquad R = \cos\gamma\cdot\dot{v} - B\sin\gamma.
\]

The equations of motion are of course,
\begin{align*}
\frac{d}{dt}\frac{dT}{dr} - \frac{dT}{dr} &= \frac{dV}{dr}, \\
\frac{d}{dt}\frac{dT}{dv} - \frac{dT}{dv} &= \frac{dV}{dv}, \\
\frac{d}{dt}\frac{dT}{d\dot{\gamma}} - \frac{dT}{d\gamma} &= \frac{dV}{d\gamma},
\end{align*}
where $V = \frac{\mu}{r} + \Omega$, if $\Omega$ is the disturbing function, taken with Lagrange's sign ($\mu = R^2$, if $-R$ is the disturbing function of the \textit{M\'{e}canique C\'{e}leste}).

To reduce these, we have in the first place
\begin{align*}
\frac{dT}{d\dot{r}} &= \dot{r} + Q = r^2\dot{Q}, \\
\frac{dT}{d\dot{v}} &= r^2\cos^2\gamma\,\dot{v} + Rr = r^2R, \\
\frac{dT}{d\dot{\gamma}} &= r^2\dot{\gamma} + Qr = r^2\dot{Q}.
\end{align*}

The equations are thus reduced to
\begin{align*}
\frac{d}{dt}(r^2\dot{Q}) - r\dot{Q}^2 + \frac{\mu}{r^2} &= \frac{dV}{dr}, \\
\frac{d}{dt}(r^2R) + r^2(Q\dot{R} + R\dot{A}\sin\gamma) &= \frac{dV}{dv}, \\
\frac{d}{dt}(r^2\dot{Q}) - r^2\cos\gamma\sin\gamma\,\dot{v}^2 + rR\dot{\gamma} &= \frac{dV}{d\gamma},
\end{align*}
and then substituting for $Q$, $R$ their values, viz.
\[
\dot{v} = \cos\gamma\cdot\dot{v} - B\sin\gamma,
\]
we find
\begin{align*}
\frac{d^2r}{dt^2} - r\dot{\gamma}^2 - r\cos^2\gamma\,\dot{v}^2 + \frac{\mu}{r^2} &= \frac{dV}{dr}, \\
\frac{d}{dt}(r^2\cos^2\gamma\,\dot{v}) + r^2(Q\dot{R} + R\dot{A}\sin\gamma) &= \frac{dV}{dv}, \\
\frac{d}{dt}(r^2\dot{\gamma}) - r^2\cos\gamma\sin\gamma\,\dot{v}^2 + rR\dot{\gamma} &= \frac{dV}{d\gamma},
\end{align*}
where
\begin{align*}
\mathfrak{A} &= r\left(-2A\frac{dv}{dt} - 2B\sin\gamma\cos\gamma\frac{d\gamma}{dt} + A^2 + B^2\sin^2\gamma\right), \\
\mathfrak{B} &= -(r^2B\sin\gamma\cos\gamma) + r^2\left(B\frac{d\gamma}{dt} - A\sin\gamma\cos\gamma\frac{dv}{dt} - AB\cos^2\gamma\right), \\
\mathfrak{C} &= (r^2A) + r^2\left(-(\cos^2\gamma - \sin^2\gamma)B\frac{dv}{dt} + A\sin\gamma\cos\gamma\right),
\end{align*}
in which
\begin{align*}
A &= \cos(v-\sigma)\sin\phi\cdot\dot{\theta} + \sin(v-\sigma)\dot{\phi}, \\
B &= \sin(v-\sigma)\sin\phi\cdot\dot{\theta} + \cos(v-\sigma)\dot{\phi},
\end{align*}
$\theta$, $\sigma$, $\phi$ being given functions of $t$ such that $d\sigma - \cos\phi\,d\theta = 0$.

The foregoing equations of motion are rigorously accurate.

\subsection*{II.}

Neglecting the terms which involve $A^2$, $AB$, $B^2$, we have
\begin{align*}
\mathfrak{A} &= 2Ar\frac{dr}{dt} + r^2\left(A\frac{d^2v}{dt^2} + \frac{dA}{dt}\frac{dv}{dt}\right) - 2r^2B\frac{d\gamma}{dt}, \\
\mathfrak{B} &= -\frac{d}{dt}(r^2B\sin\gamma\cos\gamma) + r^2\left(B\frac{d\gamma}{dt} - A\sin\gamma\cos\gamma\frac{dv}{dt}\right), \\
\mathfrak{C} &= \frac{d}{dt}(r^2A) + r^2\left(-(\cos^2\gamma - \sin^2\gamma)B\frac{dv}{dt} + A\sin\gamma\cos\gamma\right),
\end{align*}
and if we then neglect also the products of $A$ and $B$ by $\gamma$ and $\dfrac{d\gamma}{dt}$, we have
\[
\mathfrak{A} = 0, \qquad \mathfrak{B} = 0,
\]
where it may be noticed that, in order to obtain this last value of $\mathfrak{C}$, the only neglected term is a term containing $B\sin^2\gamma$.

Now, attending to the values of $A$ and $B$, we have
\[
\frac{dA}{dt} = \frac{dA}{dv}\frac{dv}{dt} + \frac{dA}{dt},
\]
where here and in the sequel $\dfrac{d}{dt}$ denotes differentiation in regard to $t$ in so far only as it enters through the quantities $\sigma$, $\theta$, $\phi$, which determine the position of the variable ecliptic.

Hence
\begin{align*}
\mathfrak{C} &= 2Ar\frac{dr}{dt} + r^2\frac{dA}{dv}\frac{dv}{dt} + r^2\frac{dA}{dt} - 2r^2B\frac{d\gamma}{dt} \\
&= 2Ar\frac{dr}{dt} + r^2\frac{dA}{dt} - 2r^2B\frac{d\gamma}{dt};
\end{align*}
and, as above,
\[
\mathfrak{A} = 0, \qquad \mathfrak{B} = 0.
\]

Let $r$, $v$, $\gamma$ be the values obtained on the supposition that $\mathfrak{C} = 0$, and $r + \delta r$, $v + \delta v$, $\gamma + \delta\gamma$, the accurate values; the first and second equations show that, neglecting the products of $\mathfrak{A}$ and $\mathfrak{B}$ into $\delta\gamma$ and $\dfrac{d\delta\gamma}{dt}$, we have $\delta r = 0$, $\delta v = 0$; so that the values of $r$ and $v$ are not affected by the variation of the ecliptic.

And then, substituting in the third equation $\gamma + \delta\gamma$ in the place of $\gamma$, and for $\cos(\gamma + \delta\gamma)$, $\sin(\gamma + \delta\gamma)$, $= \cos\gamma\sin\gamma + (\cos^2\gamma - \sin^2\gamma)\,\delta\gamma$,
\[
\frac{d^2\delta\gamma}{dt^2} + n^2\delta\gamma = \mathfrak{C},
\]
where $n^2 = \dfrac{\mu}{r^3}$; and this equation determines $\delta\gamma$.

\subsection*{III.}

To develop the value of $\mathfrak{C}$, we have
\begin{align*}
A &= \cos(v-\sigma)\sin\phi\cdot\dot{\theta} + \sin(v-\sigma)\dot{\phi}, \\
B &= \sin(v-\sigma)\sin\phi\cdot\dot{\theta} + \cos(v-\sigma)\dot{\phi},
\end{align*}
whence
\begin{align*}
\frac{dA}{dv} &= -B\sin\phi\cdot\dot{\theta} + A\dot{\phi}, \\
\frac{dB}{dv} &= A\sin\phi\cdot\dot{\theta} - B\dot{\phi},
\end{align*}
and thence
\begin{align*}
\mathfrak{C} &= 2Ar\frac{dr}{dt} + r^2\frac{dA}{dt} - 2r^2B\frac{d\gamma}{dt} \\
&= r^2\frac{dA}{dt} + 2Ar\frac{dr}{dt} - 2r^2B\frac{d\gamma}{dt}.
\end{align*}

\subsection*{IV.}

If in the disturbing function we write $s = \tan\gamma$, and attend only to the terms which involve $A$ and $B$, or (what is the same thing) the terms which depend on the variation of the ecliptic, then, neglecting $\gamma^2$ and $\gamma\dfrac{d\gamma}{dt}$, we have
\[
\frac{d\Omega}{d\gamma} = \mathfrak{C} = r^2\frac{dA}{dt}.
\]

\subsection*{V.}

The equations of motion, when the variation of the ecliptic is neglected, are
\begin{align*}
\frac{d^2r}{dt^2} - r\cos^2\gamma\left(\frac{dv}{dt}\right)^2 + \frac{\mu}{r^2} &= \frac{d\Omega}{dr}, \\
\frac{d}{dt}\left(r^2\cos^2\gamma\frac{dv}{dt}\right) &= \frac{d\Omega}{dv}, \\
\frac{d^2\gamma}{dt^2} + \frac{\mu}{r^3}\gamma &= \frac{1}{r^2}\frac{d\Omega}{d\gamma}.
\end{align*}

\subsection*{VI.}

For the further development of the disturbing function, writing $u = \dfrac{1}{r}$, and taking $\Omega$ to denote the disturbing function of the \textit{M\'{e}canique C\'{e}leste}, we have
\[
\Omega = \mu'm'u'^3\left(\frac{1}{\Delta} - u\,r'^2\cos\gamma\right),
\]
where
\[
\Delta^2 = 1 + u^2r'^2 - 2ur'\cos\gamma\cos(v-v') + u^2r'^2\sin^2\gamma\sin^2(v-v').
\]

Writing $s = \tan\gamma$, and neglecting $s^2$, we have
\[
\Omega = \mu'm'u'^3\left(\frac{1}{\Delta} - ur'^2\right),
\]
where
\[
\Delta^2 = 1 + u^2r'^2 - 2ur'\cos(v-v') + u^2r'^2s^2\sin^2(v-v').
\]

\subsection*{VII.}

If the latitude of the disturbing body is neglected, then $s' = 0$, and we have
\[
\Omega = \mu'm'u'^3\left(\frac{1}{\Delta} - ur'^2\right),
\]
where
\[
\Delta^2 = 1 + u^2r'^2 - 2ur'\cos(v-v').
\]

\subsection*{VIII.}

To reduce the equation
\[
\frac{d^2\gamma}{dt^2} + n^2\gamma = \frac{1}{r^2}\frac{d\Omega}{d\gamma},
\]
where $n^2 = \dfrac{\mu}{r^3}$, we write
\[
\gamma = \delta v\cdot\sin(v-\varpi) + \delta e\cdot\cos(v-\varpi),
\]
where $\varpi$ is the longitude of pericentre; and then, substituting and integrating, we obtain the values of $\delta v$ and $\delta e$.

\subsection*{IX.}

The equation
\[
\frac{d}{dt}\left(r^2\frac{dv}{dt}\right) = \frac{d\Omega}{dv}
\]
gives, integrating,
\[
r^2\frac{dv}{dt} = C + \int\frac{d\Omega}{dv}\,dt,
\]
where $C$ is a constant of integration.

\subsection*{X.}

The developments being carried as far as $m^2$, $e^2$, we have for the value of the disturbing function
\[
\Omega = \mu'm'u'^3\left(\frac{1}{\Delta} - ur'^2\right),
\]
and for the equation of the orbit
\[
u = nt + \epsilon + 2e\sin(nt + \epsilon - \varpi) + \tfrac{5}{4}e^2\sin 2(nt + \epsilon - \varpi).
\]

\vfill
\pagebreak

%%%%%%%%%%%%%%%%%%%%%%%%%%%%%%%%%%%%%%%%%%%%%%%%%%%%%%%%%%%%%%%%%%%%%%
% PAPER 219
%%%%%%%%%%%%%%%%%%%%%%%%%%%%%%%%%%%%%%%%%%%%%%%%%%%%%%%%%%%%%%%%%%%%%%
\setcounter{section}{218}
\section{ON SOME FORMULAE RELATING TO THE VARIATION OF A PLANET'S ORBIT}
\markboth{On Formulae for the Variation of a Planet's Orbit.}{}

\begin{center}
[From the \textit{Monthly Notices of the Royal Astronomical Society}, vol.\ xxi.\ (1863), pp.~43--47.]
\end{center}

In Hansen's Memoir, ``Auseinandersetzung einer zweckm\'{a}ssigen Methode zur Berechnung der absoluten St\'{o}rungen der kleinen Planeten,'' Abhand.\ der K.\ Sachs.\ Gesell.\ t.\ V.\ (1856), are contained, \S~8, some very elegant formulae for taking account of the variation of the plane of the orbit. These, in fact, depend upon the following geometrical theorem, viz., if (in the figure) $ABC$ is a spherical triangle; $P$, a point on the side $AB$; and $PM$, $PN$, the perpendiculars let fall from $P$ on the other two sides $AC$, $CB$; then we have
\begin{align}
\cos PM\sin(BC + CM) &= \cos PN\sin BN - \tan\tfrac{1}{2}C\cos BC(\sin PM + \sin PN), \\
\cos PM\cos(BC + CM) &= \cos PN\cos BN + \tan\tfrac{1}{2}C\sin BC(\sin PM + \sin PN).
\end{align}

These equations, in fact, give
\begin{align}
\cos PM\sin CM &= \cos PN\sin CN - \tan\tfrac{1}{2}C(\sin PM + \sin PN), \\
\cos PM\cos CM &= \cos PN\cos CN;
\end{align}
the latter of which is at once seen to be true, since joining the points $C$ and $P$, the two sides are respectively equal to $\cos CP$.

To verify the former one, write $\angle PCM = C_1$, $\angle PCN = C_2$, so that $C = C_1 - C_2$. Then, since
\[
\cos CP = \cos PM\cos CM = \cos PN\cos CN,
\]
\[
\sin PM = \sin CP\sin C_1, \qquad \sin PN = \sin CP\sin C_2,
\]
the equation becomes
\[
\cos CP(\tan CM - \tan CN) = -\tan\tfrac{1}{2}C\sin CP(\sin C_1 + \sin C_2),
\]
or since $\tan CM = \tan CP\cos C_1$, $\tan CN = \tan CP\cos C_2$, this is
\[
\cos C_1 - \cos C_2 = \tan\tfrac{1}{2}C(\sin C_1 + \sin C_2),
\]
which is identically true, in virtue of the equation $C = C_1 - C_2$; and, conversely, we have the original two equations.

Suppose that $XM$ is the ecliptic, $X$ being the origin of longitudes, $DP$ the instantaneous orbit, $D$ the departure-point therein, and $P$ the planet, $DD'$ the orthogonal trajectory of the successive positions of the orbit; and writing
\begin{align*}
\Upsilon, &\quad \text{the departure of the planet}, \\
v, &\quad \text{the longitude of ditto}, \\
y, &\quad \text{the latitude of ditto}, \\
\theta, &\quad \text{the longitude of node}, \\
\sigma, &\quad \text{the departure of ditto}, \\
\phi, &\quad \text{the inclination},
\end{align*}
then, in the figure, $DP = \Upsilon$, $XM = v$, $PM = y$, $XA = \theta$, $DA = \sigma$; $\angle A = \phi$.

The quantities $\theta$, $\sigma$, $\phi$ might be considered as altogether arbitrary; but to fix the ideas it is better to assume at once that they denote $\Theta$, the longitude of node, $\sigma_0$, the departure, $\Phi$, the inclination, for the initial position of the orbit, viz., in the figure $XB = \Theta$, $D'B = \sigma_0$, $\angle B = \Phi$.

Take $DB = \sigma_0$, $\angle B = \Phi$, $BX = \Theta$, this determines a travelling orbit of reference $XN$, and origin of longitudes $X$ therein; such that, with respect to this travelling orbit, the position of the planet's orbit is determined by $\theta$, the longitude of node, $\sigma_0$ the departure of node, $\phi$, the inclination.

We have in the triangle $ABC$, $AB = \sigma - \sigma_0$, $\angle B = \Phi$, $\angle A = 180^\circ - \phi$; and if the other parts of the triangle are represented by $BC = \omega$, $AC = \psi$, then $\omega$, $\psi$ are given in terms of $\sigma - \sigma_0$, $\Phi$, $\phi$; and we have, moreover, $XC = \theta + AC = \theta + \psi + \chi$, $XC = \sigma_0 + \omega$; that is, the position of the travelling orbit $XN$, and origin of longitudes $X$ therein, are determined by $\theta + \omega + \psi$, the longitude of node, $\sigma_0 + \omega$, the departure of node, $\Phi$, the inclination.

Suppose that in reference to this travelling orbit and origin of longitudes therein, we have $v_1$, the longitude of planet, $y_1$, the latitude of ditto, viz., in the figure $XN = v_1$ (and therefore $BN = v_1$), $PN = y_1$.

Moreover, $BC + CM = BC + AM - AC = \omega + (v - \theta) - (\theta - \Theta + \omega + \psi) = v - \Theta - \psi$, hence the two formulae become
\begin{align}
\cos y\sin(v - \Theta - \psi) &= \cos y_1\sin(v_1 - \sigma_0 - \omega) - \tan\tfrac{1}{2}\Phi\cos\omega(\sin y + \sin y_1), \\
\cos y\cos(v - \Theta - \psi) &= \cos y_1\cos(v_1 - \sigma_0 - \omega) + \tan\tfrac{1}{2}\Phi\sin\omega(\sin y + \sin y_1).
\end{align}

\vfill
\pagebreak

%%%%%%%%%%%%%%%%%%%%%%%%%%%%%%%%%%%%%%%%%%%%%%%%%%%%%%%%%%%%%%%%%%%%%%
% PAPER 220
%%%%%%%%%%%%%%%%%%%%%%%%%%%%%%%%%%%%%%%%%%%%%%%%%%%%%%%%%%%%%%%%%%%%%%
\setcounter{section}{219}
\section{NOTE ON A THEOREM OF JACOBI'S, IN RELATION TO THE PROBLEM OF THREE BODIES}
\markboth{Note on a Theorem of Jacobi's.}{}

\begin{center}
[From the \textit{Monthly Notices of the Royal Astronomical Society}, vol.\ xxii.\ (1861), pp.~76--78.]
\end{center}

The following theorem of Jacobi's (Comptes Rendus, t.~III., p.~61 (1836)) has not, I think, found its way in an explicit form into any treatise of physical astronomy.

The theorem is as follows, viz.

``Consider the movement of a point without mass round the Sun, disturbed by a planet the orbit of which is circular. Let $xyz$ be the rectangular coordinates of the disturbed body, the orbit of the disturbing planet being taken as the plane of $xy$, and the Sun as the centre of coordinates; let $a$ be the distance of the disturbing planet, $nt$ its longitude, $m$ its mass, $M$ the mass of the Sun: then we have, rigorously,
\begin{multline*}
\frac{dy}{dt}\frac{dx}{dt} - \frac{dx}{dt}\frac{dy}{dt} + M\frac{x\,dy - y\,dx}{x^2+y^2+z^2} \\
- \frac{m}{a^2}\frac{(x\,dy - y\,dx)(x\cos nt + y\sin nt)}{(x^2+y^2+z^2)^{\frac{3}{2}}\{x^2+y^2+z^2 - 2a(x\cos nt + y\sin nt) + a^2\}^{\frac{1}{2}}} \\
= \frac{m}{a^2}(x\cos nt + y\sin nt).''
\end{multline*}

This is therefore a new integral equation, which, in the problem of three bodies, subsists, as regards the terms independent of the eccentricity of the disturbing planet, and which is rigorous as regards all the powers of the mass of such planet.

In the Lunar Theory the Earth must be substituted in the place of the Sun, and the Sun taken as the disturbing planet.

To prove the theorem, as expressed in polar coordinates, I take the equations of motion in the form in which I have employed them in my ``Memoir on the Theory of Disturbed Elliptic Motion'' (Memoirs, vol.\ xxvii.\ p.~1 (1859)), [212], viz.
\begin{align}
\frac{d}{dt}\left(\frac{dr}{dt}\right) - r\left(\frac{dv}{dt}\right)^2\cos^2\gamma - r\left(\frac{d\gamma}{dt}\right)^2 + \frac{\mu}{r^2} &= \frac{d\Omega}{dr}, \\
\frac{d}{dt}\left(r^2\cos^2\gamma\frac{dv}{dt}\right) &= \frac{d\Omega}{dv}, \\
\frac{d}{dt}\left(r^2\frac{d\gamma}{dt}\right) + r^2\cos\gamma\sin\gamma\left(\frac{dv}{dt}\right)^2 &= \frac{d\Omega}{d\gamma},
\end{align}
where
\[
\Omega = \frac{m'}{r'} - m'\frac{x\xi + y\eta + z\zeta}{r'^3},
\]
or, since $\cos H = \cos\gamma\cos(v-v')$, and the Sun is considered as moving in a circular orbit (i.e.\ $r' = a$, $v' = nt$), we have
\[
\Omega = \frac{m}{\sqrt{r^2 - 2ra\cos\gamma\cos(v-nt) + a^2}},
\]
so that $\Omega$ is a function of $r$, $v$, $\gamma$ and of $t$, which last quantity enters only in the combination $v-nt$.

Hence the complete differential coefficient of $\Omega$ is
\[
\frac{d\Omega}{dt} = \frac{d\Omega}{dr}\frac{dr}{dt} + \frac{d\Omega}{dv}\frac{dv}{dt} + \frac{d\Omega}{d\gamma}\frac{d\gamma}{dt} + \frac{\partial\Omega}{\partial t},
\]
where $\dfrac{\partial\Omega}{\partial t}$ denotes, as usual, the differential coefficient in regard to the time, in so far as it enters through the coordinates $r$, $v$, $\gamma$ of the disturbed body.

We have, as usual,
\[
\frac{dr^2 + r^2(\cos^2\gamma\,dv^2 + d\gamma^2)}{dt^2} - 2\frac{d\Omega}{dt} = -2\frac{d\Omega}{dt},
\]
and, from the foregoing equation,
\[
\frac{dv}{dt}\frac{d\Omega}{dv} = \frac{dv}{dt}\frac{d}{dt}\left(r^2\cos^2\gamma\frac{dv}{dt}\right);
\]
hence, substituting this value, transposing, and integrating, we have
\[
r^2\cos^2\gamma\frac{dv}{dt} + \int\frac{d\Omega}{dt}\,dt = \text{const.},
\]
which is Jacobi's equation expressed in terms of the coordinates $r$, $v$, $\gamma$.

M.\ de Pont\'{e}coulant, in his Lunar Theory (1846), where the solar eccentricity is neglected, writes (p.~91),
\[
\frac{r^2\,dv}{(1+s^2)\,dt} = m'h - \int dR,
\]
where
\[
r^2\frac{dv}{dt} = mn = m, \quad \text{since $n$ is there put equal to unity};
\]
and combining this with the equation (p.~43),
\[
r^2\frac{dv}{dt} - mh = \int\frac{dR}{dv}\,dv,
\]
we have
\[
\int dR = R - mh + \frac{mr^2\,dv}{(1+s^2)\,dt},
\]
and substituting this value of $\int dR$ in the integral of Vis Viva (p.~41),
\[
\tfrac{1}{2}\left(\frac{dr^2+r^2(dv^2+d\gamma^2)}{dt^2}\right) - \frac{M}{r} = \Omega,
\]
we have what is, in fact, Jacobi's equation.

\vfill
\pagebreak

%%%%%%%%%%%%%%%%%%%%%%%%%%%%%%%%%%%%%%%%%%%%%%%%%%%%%%%%%%%%%%%%%%%%%%
% PAPER 221
%%%%%%%%%%%%%%%%%%%%%%%%%%%%%%%%%%%%%%%%%%%%%%%%%%%%%%%%%%%%%%%%%%%%%%
\setcounter{section}{220}
\section{ON THE SECULAR ACCELERATION OF THE MOON'S MEAN MOTION}
\markboth{On the Secular Acceleration of the Moon's Mean Motion.}{}

\begin{center}
[From the \textit{Monthly Notices of the Royal Astronomical Society}, vol.\ xxii.\ (1862), pp.~171--231.]
\end{center}

The present Memoir exhibits a new method of taking account, in the Lunar Theory, of the variation of the eccentricity of the Sun's orbit.

The approximation is carried to the same extent as in Prof.\ Adams' Memoir ``On the Secular Variation of the Moon's Mean Motion'' (Phil.\ Trans., vol.\ CXLIII.\ (1853), pp.~397--406); and I obtain results agreeing precisely with his, viz., besides his new periodic terms in the longitude and radius vector, I obtain in the longitude the secular term
\[
\left(\tfrac{3}{4}m^2 - \tfrac{81}{32}m^4\right)\int\frac{a}{a'}(e'^2 - E'^2)\,n\,dt,
\]
and in the quotient radius, or radius vector divided by the mean distance, the secular term
\[
\left(\tfrac{3}{8}m^2 - \tfrac{81}{64}m^4\right)\frac{a}{a'}(e'^2 - E'^2)\,nt,
\]
which is, in fact, as will be shown, included implicitly in the results given in Professor Adams' Memoir.

In quoting the foregoing results, I have written $e'^2 - E'^2$ in the place of $(e' + f't)^2 - e'^2 = 2e'f't$, which in the notation of the present Memoir it should have been; and I purposely refrain from here explaining the precise signification of the symbols: this is carefully done in the sequel.

The method appears to me a very simple one in principle; and it possesses the advantage that it is not incorporated step by step with a lunar theory in which the eccentricity of the Sun's orbit is treated as constant; but it is added on to such a lunar theory, giving in the Moon's coordinates the supplementary terms which arise from the variation of the solar eccentricity, and thus serving as a verification of any process employed for taking account of such variation.

I have given the details of the work in a series of Annexes, 1 to 23: this appears to me the best course for presenting the investigation in a readable form.

\subsection*{I.}

The inclination and eccentricity of the Moon's orbit, and, \textit{a fortiori}, the variation of the position of the Ecliptic, and the Sun's latitude, are neglected; and the longitudes are measured from a fixed point in the Ecliptic.

I write $n$, the actual mean motion of the Moon at a given epoch; viz., it is assumed that the mean longitude at the time $t$ is $\epsilon + nt + n_2t^2 + \cdots$ where $\epsilon$, $n$, $n_2$, etc.\ are absolute constants; and, moreover, $a$, the calculated mean distance of the Moon; that is, $n^2a^3$ is the sum of the masses of the Earth and Moon; $a$ is therefore an absolute constant; and, in like manner, $n'$, the actual mean motion of the Sun at the same epoch, $a'$, the calculated mean distance of the Sun; that is, if it were necessary to pay attention to the secular variation of the mean motion of the Sun, the assumption would be that the mean longitude was $\epsilon' + n't + n'_2t^2 + \cdots$, $\epsilon'$, $n'$, $n'_2$, etc.\ being absolute constants, and $n'^2a'^3$ the sum of the masses of the Sun and Earth; $a'$ would thus also be an absolute constant.

But for the purpose of the present investigation the secular variation of the mean longitude of the Sun is neglected, or it is assumed that the mean longitude of the Sun is $\epsilon' + n't$, $\epsilon'$, $n'$ being absolute constants; and that $n'^2a'^3$ is the sum of the masses of the Sun and Earth, $a'$ being thus also an absolute constant.

I put also $m$, the ratio of the mean motions of the Sun and Moon; that is,
\[
m = \frac{n'}{n}, \quad \text{or } n' = mn;
\]
$m$ is also an absolute constant.

The Sun is considered as moving in an elliptic orbit, the eccentricity whereof is $e' + \delta e'$ or $e' + f't$, $e'$ and $f'$ being absolute constants; the longitude of the Sun's perigee may be taken to be $\varpi' + (1-c')n't$; so that the mean anomaly $g$ is $= n't + \epsilon' - [\varpi' + (1-c')n't] = c'n't + \epsilon' - \varpi'$; $c'$, $\epsilon'$, $\varpi'$ being absolute constants; but $c'$ is in fact treated as being $= 1$.

Hence, if $r'$, $v'$ are the radius vector and longitude of the Sun, we have
\begin{align*}
r' &= a'\,\mathrm{elqr}(e' + \delta e', g), \\
v' &= \varpi' + (1-c')n't + \mathrm{elta}(e' + \delta e', g) \\
&= n't + \epsilon' + [\mathrm{elta}(e' + \delta e', g) - g],
\end{align*}
where $g = c'n't + \text{const}$.

In the expression for the disturbing function the Sun's mass is taken to be $= n'^2a'^3$, or, what is the same thing, $= m^2n^2a'^3$.

Let $r$, $v$ be the radius vector and longitude of the Moon; then, taking the usual approximate expression of the Disturbing Function, the equations of motion are
\begin{align}
\frac{d}{dt}\left(\frac{dr}{dt}\right) - r\left(\frac{dv}{dt}\right)^2 + \frac{n^2a^3}{r^2} &= \frac{dR}{dr}, \\
\frac{d}{dt}\left(r^2\frac{dv}{dt}\right) &= \frac{dR}{dv}.
\end{align}

It will be convenient to assume $\rho$, the quotient radius of the Moon's orbit, $\rho'$, the quotient radius of the Sun's orbit; that is
\[
r = \rho a, \qquad r' = \rho' a'.
\]

The equations of motion thus become
\begin{align}
\frac{d}{dt}\left(\frac{d\rho}{dt}\right) - \rho\left(\frac{dv}{dt}\right)^2 + \frac{n^2}{\rho^2} &= m^2n^2\,\frac{dP}{d\rho}, \\
\frac{d}{dt}\left(\rho^2\frac{dv}{dt}\right) &= m^2n^2\,\frac{dP}{dv},
\end{align}
where for shortness
\[
P = \frac{\rho}{\rho'^3}\left(\tfrac{3}{2}\sin^2(v-v') - \tfrac{1}{2}\right),
\]
in which
\[
\rho' = \mathrm{elqr}(e' + \delta e', g), \quad v' = n't + \epsilon' + [\mathrm{elta}(e', g) - g].
\]

I now change the notation by writing $\rho + \delta\rho$, $v + \delta v$, in the place of $\rho$, $v$, respectively, using henceforward $\rho$, $v$ to denote
\[
\rho = \mathrm{elqr}(e', g), \quad v = n't + \epsilon' + [\mathrm{elta}(e', g) - g];
\]
and I write also $\rho + \delta\rho$, $v + \delta v$, in the place of $\rho$, $v$, using henceforward $\rho$, $v$ to denote the solutions of the equations obtained from the equations of motion by writing therein $\rho$, $v$ instead of the complete values $\rho + \delta\rho$, $v + \delta v$.

Suppose, in like manner, that the complete values of $P$, $Q$ are denoted by $P + \delta P$, $Q + \delta Q$, where
\[
\delta P = \frac{dP}{d\rho}\delta\rho + \frac{dP}{dv}\delta v + \frac{dP}{d\rho'}\delta\rho' + \frac{dP}{dv'}\delta v',
\]
with a like value for $\delta Q$, the first powers of $\delta\rho$, $\delta v$, $\delta\rho'$, $\delta v'$ being alone attended to.

Then, observing that the equations of motion are satisfied when $\delta\rho$, $\delta v$, $\delta\rho'$, $\delta v'$ are neglected, we have
\begin{align}
\frac{d}{dt}\left(\frac{d\delta\rho}{dt}\right) - \left(\frac{dv}{dt}\right)^2\delta\rho - 2\rho\frac{dv}{dt}\frac{d\delta v}{dt} - \frac{2n^2}{\rho^3}\delta\rho &= m^2n^2\,\delta P, \\
\frac{d}{dt}\left(\rho^2\frac{d\delta v}{dt} + 2\rho\frac{dv}{dt}\delta\rho\right) &= m^2n^2\,\delta Q.
\end{align}

The second of these equations gives
\[
\rho^2\frac{d\delta v}{dt} + 2\rho\frac{dv}{dt}\delta\rho = C + m^2n^2\int\delta Q\,dt,
\]
where the constant of integration, $C$, is to be so determined that $\delta v$ may not contain any term of the form $kt$ (for any such term is taken to be included in the term $nt$ of $v + \delta v$).

Multiplying the equation just obtained by $-\dfrac{2}{\rho}\dfrac{dv}{dt}$, and adding it to the first equation, we have
\begin{multline}
\frac{d^2\delta\rho}{dt^2} + \left(\frac{2n^2}{\rho^3} - \left(\frac{dv}{dt}\right)^2\right)\delta\rho - \frac{2}{\rho}\frac{dv}{dt}\left(C + m^2n^2\int\delta Q\,dt\right) \\
= m^2n^2\left(\delta P + \frac{2}{\rho}\frac{dv}{dt}\delta Q\right),
\end{multline}
which, with the above-mentioned integral equation, are the equations for the solution of the problem; but it will be convenient to write them under the slightly different form
\begin{align}
\frac{d^2\delta\rho}{dt^2} + n^2\delta\rho &= \left(n^2 + \frac{2n^2}{\rho^3} - \left(\frac{dv}{dt}\right)^2\right)\delta\rho + m^2n^2\left(\delta P + \frac{2}{\rho}\frac{dv}{dt}(C + \int\delta Q\,dt)\right), \\
\frac{d\delta v}{dt} &= -\frac{2}{\rho}\frac{dv}{dt}\delta\rho + \frac{1}{\rho^2}\left(C + m^2n^2\int\delta Q\,dt\right).
\end{align}

In these equations $C$ is determined, as above, by the condition that $\dfrac{d\delta v}{dt}$ may contain no constant term; the values of $\rho'$, $v'$, $\delta\rho'$, $\delta v'$ are of course given by the theory of elliptic motion, and those of $\rho$, $v$ are given by the ordinary lunar theory, in which the eccentricity of the solar orbit is treated as a constant; and then, $\delta\rho$, $\delta v$ being obtained by integrating the equations, the radius vector and longitude of the Moon are $a(\rho + \delta\rho)$ and $v + \delta v$ respectively.

We have
\[
Q = \frac{\rho}{\rho'^3}\left(-\tfrac{3}{2}\sin 2(v-v') - 2(v-v')\right).
\]

Moreover, by the lunar theory, observing that Plana's $a$ is, or may be considered, identical with the $a$ of the present Memoir, and putting also
\[
2\nu = (2-2m)nt + \text{const.}, \qquad g = c'mn't + \text{const.},
\]
we have
\begin{align}
\frac{1}{\rho} &= 1 + \tfrac{1}{6}m^2 - \tfrac{1}{3}m^2e'^2 - \tfrac{1}{4}m^2e\cos g + \tfrac{1}{8}m^2e^2\cos 2g \nonumber\\
&\quad + \tfrac{3}{8}m^2e'^2\cos 2\nu + \cdots, \\
v &= nt + \epsilon - \tfrac{3}{4}me\sin g + \tfrac{1}{2}m^2\sin 2\nu + \tfrac{11}{8}m^2e\sin(g+2\nu) \nonumber\\
&\quad - \tfrac{55}{24}m^2e\sin(2\nu-g) + \cdots,
\end{align}
where the series are carried as far as $m^2$ and $e^2$; the terms in $e'^2$ are given, as I shall have occasion to refer to them, but they are not used in the investigation, and, omitting them, the values are
\begin{align}
\frac{1}{\rho} &= 1 + \tfrac{1}{6}m^2 - \tfrac{1}{4}m^2e\cos g + \tfrac{1}{8}m^2e^2\cos 2g + \tfrac{3}{8}m^2e'^2\cos 2\nu, \\
v &= nt + \epsilon - \tfrac{3}{4}me\sin g + \tfrac{1}{2}m^2\sin 2\nu + \tfrac{11}{8}m^2e\sin(g+2\nu) \nonumber\\
&\quad - \tfrac{55}{24}m^2e\sin(2\nu-g).
\end{align}

For the coordinates of the Sun we have
\begin{align}
\frac{a'}{r'} &= 1 + e'\cos g + e'^2\cos 2g, \\
v' &= n't + \epsilon' + 2e'\sin g + \tfrac{5}{4}e'^2\sin 2g,
\end{align}
the series being carried as far as $e'^2$; but the terms in $e'^2$ are only used for the formation of $\delta\rho'$, $\delta v'$; and, omitting them, we have
\[
\frac{a'}{r'} = 1 + e'\cos g, \qquad v' = n't + \epsilon' + 2e'\sin g.
\]

If $e' + f't$ is written for $e'$, then the value of $\delta e'$ is $= f't$; but as only the terms multiplied by the simple power $f'$ are attended to, we may for convenience write $\delta e' = t$, the factor $f'$ being restored in the final results: we thus have
\begin{align*}
\delta\rho' &= -t\cos g + 2e't\cos 2g, \\
\delta v' &= -2t\sin g + \tfrac{5}{2}e't\sin 2g,
\end{align*}
and we may add the equations
\begin{align*}
\frac{dv'}{dt} &= n' + 2e'n'\cos g + \tfrac{5}{4}e'^2n'\cos 2g, \\
\delta\rho' &= -3e't - 3t\cos g + \tfrac{5}{2}e't\cos 2g, \\
\delta v' &= -2t\sin g - \tfrac{5}{4}e't\sin 2g,
\end{align*}
which will be found useful.

\subsection*{II.}

Proceeding now to the development of the solution, we have
\begin{multline}
\delta P = \frac{1}{\rho'^3}\left(1 + \tfrac{3}{2}\cos 2(v-v')\right)\delta\rho \\
+ \frac{3}{\rho'^3}\left(\tfrac{1}{2}\sin 2(v-v')\right)\delta v + \frac{3}{2\rho'^3}\sin 2(v-v')\frac{\delta\rho'}{\rho'} + \cdots,
\end{multline}
where the terms containing $\delta\rho$ and $\delta v$ are (see Annex 1)
\begin{multline*}
\left(1 + \tfrac{3}{2}e'^2\right)\delta\rho + \tfrac{3}{2}e'\cos\nu\,\delta v + \tfrac{3}{2}e'\sin\nu\,\frac{\delta\rho'}{\rho'} + \cdots \\
= \left(1 + \tfrac{3}{2}e'^2\right)\left(-\tfrac{3}{4}\sin 2\nu - \tfrac{1}{2}e\sin(2\nu-g) + \tfrac{3}{4}e\sin(2\nu+g)\right) \\
+ \tfrac{3}{2}e'\cos\nu\left(\tfrac{1}{2}\cos 2\nu + e\cos(2\nu-g) - e\cos(2\nu+g)\right) + \cdots
\end{multline*}
and also
\begin{multline*}
\tfrac{3}{2}\sin 2\nu\left(-\tfrac{3}{4}\sin 2\nu + \cdots\right) + \tfrac{3}{2}\cos 2\nu\left(\tfrac{1}{2}\cos 2\nu + \cdots\right) \\
= \tfrac{3}{4}\left(-\tfrac{3}{2}\sin^2 2\nu + \cos^2 2\nu\right) + \cdots
\end{multline*}
where the terms containing $\delta\rho'$ and $\delta v'$ are (see Annex 2)
\begin{multline*}
\tfrac{3}{2}\sin 2\nu\left(-t\cos g + 2e't\cos 2g\right) \\
+ \tfrac{3}{2}\cos 2\nu\left(-2t\sin g + \tfrac{5}{2}e't\sin 2g\right) + \cdots
\end{multline*}

The development being continued through the annexes referred to, the final results are obtained for the secular acceleration of the Moon's mean motion, confirming those of Professor Adams.

\end{document}
